% Options for packages loaded elsewhere
\PassOptionsToPackage{unicode}{hyperref}
\PassOptionsToPackage{hyphens}{url}
%
\documentclass[
]{article}
\usepackage{amsmath,amssymb}
\usepackage{lmodern}
\usepackage{iftex}
\ifPDFTeX
  \usepackage[T1]{fontenc}
  \usepackage[utf8]{inputenc}
  \usepackage{textcomp} % provide euro and other symbols
\else % if luatex or xetex
  \usepackage{unicode-math}
  \defaultfontfeatures{Scale=MatchLowercase}
  \defaultfontfeatures[\rmfamily]{Ligatures=TeX,Scale=1}
\fi
% Use upquote if available, for straight quotes in verbatim environments
\IfFileExists{upquote.sty}{\usepackage{upquote}}{}
\IfFileExists{microtype.sty}{% use microtype if available
  \usepackage[]{microtype}
  \UseMicrotypeSet[protrusion]{basicmath} % disable protrusion for tt fonts
}{}
\makeatletter
\@ifundefined{KOMAClassName}{% if non-KOMA class
  \IfFileExists{parskip.sty}{%
    \usepackage{parskip}
  }{% else
    \setlength{\parindent}{0pt}
    \setlength{\parskip}{6pt plus 2pt minus 1pt}}
}{% if KOMA class
  \KOMAoptions{parskip=half}}
\makeatother
\usepackage{xcolor}
\usepackage[margin=1in]{geometry}
\usepackage{color}
\usepackage{fancyvrb}
\newcommand{\VerbBar}{|}
\newcommand{\VERB}{\Verb[commandchars=\\\{\}]}
\DefineVerbatimEnvironment{Highlighting}{Verbatim}{commandchars=\\\{\}}
% Add ',fontsize=\small' for more characters per line
\usepackage{framed}
\definecolor{shadecolor}{RGB}{248,248,248}
\newenvironment{Shaded}{\begin{snugshade}}{\end{snugshade}}
\newcommand{\AlertTok}[1]{\textcolor[rgb]{0.94,0.16,0.16}{#1}}
\newcommand{\AnnotationTok}[1]{\textcolor[rgb]{0.56,0.35,0.01}{\textbf{\textit{#1}}}}
\newcommand{\AttributeTok}[1]{\textcolor[rgb]{0.77,0.63,0.00}{#1}}
\newcommand{\BaseNTok}[1]{\textcolor[rgb]{0.00,0.00,0.81}{#1}}
\newcommand{\BuiltInTok}[1]{#1}
\newcommand{\CharTok}[1]{\textcolor[rgb]{0.31,0.60,0.02}{#1}}
\newcommand{\CommentTok}[1]{\textcolor[rgb]{0.56,0.35,0.01}{\textit{#1}}}
\newcommand{\CommentVarTok}[1]{\textcolor[rgb]{0.56,0.35,0.01}{\textbf{\textit{#1}}}}
\newcommand{\ConstantTok}[1]{\textcolor[rgb]{0.00,0.00,0.00}{#1}}
\newcommand{\ControlFlowTok}[1]{\textcolor[rgb]{0.13,0.29,0.53}{\textbf{#1}}}
\newcommand{\DataTypeTok}[1]{\textcolor[rgb]{0.13,0.29,0.53}{#1}}
\newcommand{\DecValTok}[1]{\textcolor[rgb]{0.00,0.00,0.81}{#1}}
\newcommand{\DocumentationTok}[1]{\textcolor[rgb]{0.56,0.35,0.01}{\textbf{\textit{#1}}}}
\newcommand{\ErrorTok}[1]{\textcolor[rgb]{0.64,0.00,0.00}{\textbf{#1}}}
\newcommand{\ExtensionTok}[1]{#1}
\newcommand{\FloatTok}[1]{\textcolor[rgb]{0.00,0.00,0.81}{#1}}
\newcommand{\FunctionTok}[1]{\textcolor[rgb]{0.00,0.00,0.00}{#1}}
\newcommand{\ImportTok}[1]{#1}
\newcommand{\InformationTok}[1]{\textcolor[rgb]{0.56,0.35,0.01}{\textbf{\textit{#1}}}}
\newcommand{\KeywordTok}[1]{\textcolor[rgb]{0.13,0.29,0.53}{\textbf{#1}}}
\newcommand{\NormalTok}[1]{#1}
\newcommand{\OperatorTok}[1]{\textcolor[rgb]{0.81,0.36,0.00}{\textbf{#1}}}
\newcommand{\OtherTok}[1]{\textcolor[rgb]{0.56,0.35,0.01}{#1}}
\newcommand{\PreprocessorTok}[1]{\textcolor[rgb]{0.56,0.35,0.01}{\textit{#1}}}
\newcommand{\RegionMarkerTok}[1]{#1}
\newcommand{\SpecialCharTok}[1]{\textcolor[rgb]{0.00,0.00,0.00}{#1}}
\newcommand{\SpecialStringTok}[1]{\textcolor[rgb]{0.31,0.60,0.02}{#1}}
\newcommand{\StringTok}[1]{\textcolor[rgb]{0.31,0.60,0.02}{#1}}
\newcommand{\VariableTok}[1]{\textcolor[rgb]{0.00,0.00,0.00}{#1}}
\newcommand{\VerbatimStringTok}[1]{\textcolor[rgb]{0.31,0.60,0.02}{#1}}
\newcommand{\WarningTok}[1]{\textcolor[rgb]{0.56,0.35,0.01}{\textbf{\textit{#1}}}}
\usepackage{graphicx}
\makeatletter
\def\maxwidth{\ifdim\Gin@nat@width>\linewidth\linewidth\else\Gin@nat@width\fi}
\def\maxheight{\ifdim\Gin@nat@height>\textheight\textheight\else\Gin@nat@height\fi}
\makeatother
% Scale images if necessary, so that they will not overflow the page
% margins by default, and it is still possible to overwrite the defaults
% using explicit options in \includegraphics[width, height, ...]{}
\setkeys{Gin}{width=\maxwidth,height=\maxheight,keepaspectratio}
% Set default figure placement to htbp
\makeatletter
\def\fps@figure{htbp}
\makeatother
\setlength{\emergencystretch}{3em} % prevent overfull lines
\providecommand{\tightlist}{%
  \setlength{\itemsep}{0pt}\setlength{\parskip}{0pt}}
\setcounter{secnumdepth}{-\maxdimen} % remove section numbering
\ifLuaTeX
  \usepackage{selnolig}  % disable illegal ligatures
\fi
\IfFileExists{bookmark.sty}{\usepackage{bookmark}}{\usepackage{hyperref}}
\IfFileExists{xurl.sty}{\usepackage{xurl}}{} % add URL line breaks if available
\urlstyle{same} % disable monospaced font for URLs
\hypersetup{
  pdftitle={Regression Model},
  pdfauthor={Wilbert},
  hidelinks,
  pdfcreator={LaTeX via pandoc}}

\title{Regression Model}
\author{Wilbert}
\date{2023-06-14}

\begin{document}
\maketitle

\hypertarget{regression-vs-classification}{%
\section{Regression vs
Classification}\label{regression-vs-classification}}

\begin{itemize}
\tightlist
\item
  Regression -\textgreater{} Continuous
\item
  Classification -\textgreater{} Discrete (Labels \& Categorical)
\end{itemize}

Regression and Classification is differentiated by it's dependent
variable.

\hypertarget{data-description}{%
\section{1. Data Description}\label{data-description}}

\url{https://archive.ics.uci.edu/dataset/9/auto+mpg}

\begin{verbatim}
1. mpg:           continuous (Miles per Gallon, in gallon)
2. cylinders:     multi-valued discrete (Number of Cylinders)
3. displacement:  continuous (Engine Capacity)
4. horsepower:    continuous (Engine Power)
5. weight:        continuous (Vehicle Weight)
6. acceleration:  continuous (Vehicle Acceleration)
7. model year:    multi-valued discrete
8. origin:        multi-valued discrete
9. car name:      string (unique for each instance)
\end{verbatim}

\hypertarget{load-libraries-data}{%
\section{2. Load libraries \& data}\label{load-libraries-data}}

\begin{Shaded}
\begin{Highlighting}[]
\FunctionTok{library}\NormalTok{(rpart)}
\FunctionTok{library}\NormalTok{(rpart.plot)}
\FunctionTok{library}\NormalTok{(MLmetrics) }
\end{Highlighting}
\end{Shaded}

\begin{verbatim}
## 
## Attaching package: 'MLmetrics'
\end{verbatim}

\begin{verbatim}
## The following object is masked from 'package:base':
## 
##     Recall
\end{verbatim}

\begin{Shaded}
\begin{Highlighting}[]
\CommentTok{\# Regression use RMSF \& Rsquared, pROC {-}\textgreater{} ROC}
\CommentTok{\# Classification Sensivity, Specificity, Accuracy}

\NormalTok{autompg }\OtherTok{\textless{}{-}} \FunctionTok{read.table}\NormalTok{(}\StringTok{"auto{-}mpg.data"}\NormalTok{)}
\FunctionTok{colnames}\NormalTok{(autompg) }\OtherTok{\textless{}{-}} \FunctionTok{c}\NormalTok{(}\StringTok{"mpg"}\NormalTok{, }\StringTok{"cylinders"}\NormalTok{, }\StringTok{"displacement"}\NormalTok{, }\StringTok{"horsepower"}\NormalTok{, }\StringTok{"weight"}\NormalTok{, }\StringTok{"acceleration"}\NormalTok{, }\StringTok{"model\_year"}\NormalTok{, }\StringTok{"origin"}\NormalTok{, }\StringTok{"car\_name"}\NormalTok{)}
\FunctionTok{head}\NormalTok{(autompg)}
\end{Highlighting}
\end{Shaded}

\begin{verbatim}
##   mpg cylinders displacement horsepower weight acceleration model_year origin
## 1  18         8          307      130.0   3504         12.0         70      1
## 2  15         8          350      165.0   3693         11.5         70      1
## 3  18         8          318      150.0   3436         11.0         70      1
## 4  16         8          304      150.0   3433         12.0         70      1
## 5  17         8          302      140.0   3449         10.5         70      1
## 6  15         8          429      198.0   4341         10.0         70      1
##                    car_name
## 1 chevrolet chevelle malibu
## 2         buick skylark 320
## 3        plymouth satellite
## 4             amc rebel sst
## 5               ford torino
## 6          ford galaxie 500
\end{verbatim}

\hypertarget{statistics-summary}{%
\section{3. Statistics Summary}\label{statistics-summary}}

\begin{Shaded}
\begin{Highlighting}[]
\NormalTok{basicSummary }\OtherTok{\textless{}{-}} \ControlFlowTok{function}\NormalTok{(df, }\AttributeTok{dgts =} \DecValTok{3}\NormalTok{) \{}
\NormalTok{  m }\OtherTok{\textless{}{-}} \FunctionTok{ncol}\NormalTok{(df) }\CommentTok{\# column}
\NormalTok{  varNames }\OtherTok{\textless{}{-}} \FunctionTok{colnames}\NormalTok{(df)}
\NormalTok{  varType }\OtherTok{\textless{}{-}} \FunctionTok{vector}\NormalTok{(}\StringTok{"character"}\NormalTok{, m)}
\NormalTok{  topLevel }\OtherTok{\textless{}{-}} \FunctionTok{vector}\NormalTok{(}\StringTok{"character"}\NormalTok{, m)}
\NormalTok{  topCount }\OtherTok{\textless{}{-}} \FunctionTok{vector}\NormalTok{(}\StringTok{"numeric"}\NormalTok{, m)}
\NormalTok{  missCount }\OtherTok{\textless{}{-}} \FunctionTok{vector}\NormalTok{(}\StringTok{"numeric"}\NormalTok{, m)}
\NormalTok{  levels }\OtherTok{\textless{}{-}} \FunctionTok{vector}\NormalTok{(}\StringTok{"numeric"}\NormalTok{, m)}
  
  \ControlFlowTok{for}\NormalTok{ (i }\ControlFlowTok{in} \DecValTok{1}\SpecialCharTok{:}\NormalTok{m) \{}
\NormalTok{    x }\OtherTok{\textless{}{-}}\NormalTok{ df[,i]}
\NormalTok{    varType[i] }\OtherTok{\textless{}{-}} \FunctionTok{class}\NormalTok{(x)}
\NormalTok{    xtab }\OtherTok{\textless{}{-}} \FunctionTok{table}\NormalTok{(x, }\AttributeTok{useNA =} \StringTok{"ifany"}\NormalTok{)}
\NormalTok{    levels[i] }\OtherTok{\textless{}{-}} \FunctionTok{length}\NormalTok{(xtab)}
\NormalTok{    nums }\OtherTok{\textless{}{-}} \FunctionTok{as.numeric}\NormalTok{(xtab)}
\NormalTok{    maxnum }\OtherTok{\textless{}{-}} \FunctionTok{max}\NormalTok{(nums)}
\NormalTok{    topCount[i] }\OtherTok{\textless{}{-}}\NormalTok{ maxnum}
\NormalTok{    maxIndex }\OtherTok{\textless{}{-}} \FunctionTok{which.max}\NormalTok{(nums)}
\NormalTok{    lvls }\OtherTok{\textless{}{-}} \FunctionTok{names}\NormalTok{(xtab)}
\NormalTok{    topLevel[i] }\OtherTok{\textless{}{-}}\NormalTok{ lvls[maxIndex]}
\NormalTok{    missIndex }\OtherTok{\textless{}{-}} \FunctionTok{which}\NormalTok{((}\FunctionTok{is.na}\NormalTok{(x)) }\SpecialCharTok{|}\NormalTok{ (x }\SpecialCharTok{==} \StringTok{""}\NormalTok{) }\SpecialCharTok{|}\NormalTok{ (x }\SpecialCharTok{==} \StringTok{" "}\NormalTok{))}
\NormalTok{    missCount[i] }\OtherTok{\textless{}{-}} \FunctionTok{length}\NormalTok{(missIndex)}
\NormalTok{  \}}
\NormalTok{  n }\OtherTok{\textless{}{-}} \FunctionTok{nrow}\NormalTok{(df) }\CommentTok{\# Rows}
\NormalTok{  topFrac }\OtherTok{\textless{}{-}} \FunctionTok{round}\NormalTok{(topCount}\SpecialCharTok{/}\NormalTok{n, }\AttributeTok{digits =}\NormalTok{ dgts)}
\NormalTok{  missFrac }\OtherTok{\textless{}{-}} \FunctionTok{round}\NormalTok{(missCount}\SpecialCharTok{/}\NormalTok{n, }\AttributeTok{digits =}\NormalTok{ dgts)}
\NormalTok{  summaryFrame }\OtherTok{\textless{}{-}} \FunctionTok{data.frame}\NormalTok{(}\AttributeTok{variable =}\NormalTok{ varNames, }
                             \AttributeTok{type =}\NormalTok{ varType, }
                             \AttributeTok{levels =}\NormalTok{ levels,}
                             \AttributeTok{topLevel =}\NormalTok{ topLevel,}
                             \AttributeTok{topCount =}\NormalTok{ topCount,}
                             \AttributeTok{topFrac =}\NormalTok{ topFrac,}
                             \AttributeTok{missFreq =}\NormalTok{ missCount,}
                             \AttributeTok{missFrac =}\NormalTok{ missFrac)}
  \FunctionTok{return}\NormalTok{(summaryFrame)}
\NormalTok{\}}
\end{Highlighting}
\end{Shaded}

\begin{Shaded}
\begin{Highlighting}[]
\FunctionTok{basicSummary}\NormalTok{(autompg)}
\end{Highlighting}
\end{Shaded}

\begin{verbatim}
##       variable      type levels   topLevel topCount topFrac missFreq missFrac
## 1          mpg   numeric    129         13       20   0.050        0        0
## 2    cylinders   integer      5          4      204   0.513        0        0
## 3 displacement   numeric     82         97       21   0.053        0        0
## 4   horsepower character     94      150.0       22   0.055        0        0
## 5       weight   numeric    351       1985        4   0.010        0        0
## 6 acceleration   numeric     95       14.5       23   0.058        0        0
## 7   model_year   integer     13         73       40   0.101        0        0
## 8       origin   integer      3          1      249   0.626        0        0
## 9     car_name character    305 ford pinto        6   0.015        0        0
\end{verbatim}

\begin{Shaded}
\begin{Highlighting}[]
\CommentTok{\# Because Horse wpoer is not the right data type, we must change it to the correct data type}
\end{Highlighting}
\end{Shaded}

\hypertarget{data-preparation}{%
\section{4. Data Preparation}\label{data-preparation}}

In Summary : Data Problems.

\begin{enumerate}
\def\labelenumi{\arabic{enumi}.}
\tightlist
\item
  car\_name is meaningless -\textgreater{} ``remove''
\item
  convert all var into numerical
\item
  Horsepower have missng value
\end{enumerate}

Rule of thumb : if missing value \textless{} 20\% of total data then we
can remove it, if not then median is recomended

\begin{Shaded}
\begin{Highlighting}[]
\CommentTok{\# Remove car\_name}
\NormalTok{autompg }\OtherTok{\textless{}{-}} \FunctionTok{subset}\NormalTok{(autompg, }\AttributeTok{select =} \SpecialCharTok{{-}}\FunctionTok{c}\NormalTok{(car\_name))}

\CommentTok{\# Change all var into numerical}
\NormalTok{autompg}\SpecialCharTok{$}\NormalTok{horsepower }\OtherTok{\textless{}{-}} \FunctionTok{as.numeric}\NormalTok{(autompg}\SpecialCharTok{$}\NormalTok{horsepower)}
\end{Highlighting}
\end{Shaded}

\begin{verbatim}
## Warning: NAs introduced by coercion
\end{verbatim}

\begin{Shaded}
\begin{Highlighting}[]
\NormalTok{autompg}\SpecialCharTok{$}\NormalTok{cylinders }\OtherTok{\textless{}{-}} \FunctionTok{as.numeric}\NormalTok{(autompg}\SpecialCharTok{$}\NormalTok{cylinders)}
\NormalTok{autompg}\SpecialCharTok{$}\NormalTok{model\_year }\OtherTok{\textless{}{-}} \FunctionTok{as.numeric}\NormalTok{(autompg}\SpecialCharTok{$}\NormalTok{model\_year)}
\NormalTok{autompg}\SpecialCharTok{$}\NormalTok{origin }\OtherTok{\textless{}{-}} \FunctionTok{as.numeric}\NormalTok{(autompg}\SpecialCharTok{$}\NormalTok{origin)}

\CommentTok{\# Change missing value with mean}
\NormalTok{autompg}\SpecialCharTok{$}\NormalTok{horsepower[}\FunctionTok{is.na}\NormalTok{(autompg}\SpecialCharTok{$}\NormalTok{horsepower)] }\OtherTok{\textless{}{-}} \FunctionTok{mean}\NormalTok{(autompg}\SpecialCharTok{$}\NormalTok{horsepower, }\AttributeTok{na.rm =}\NormalTok{ T)}
\FunctionTok{basicSummary}\NormalTok{(autompg)}
\end{Highlighting}
\end{Shaded}

\begin{verbatim}
##       variable    type levels topLevel topCount topFrac missFreq missFrac
## 1          mpg numeric    129       13       20   0.050        0        0
## 2    cylinders numeric      5        4      204   0.513        0        0
## 3 displacement numeric     82       97       21   0.053        0        0
## 4   horsepower numeric     94      150       22   0.055        0        0
## 5       weight numeric    351     1985        4   0.010        0        0
## 6 acceleration numeric     95     14.5       23   0.058        0        0
## 7   model_year numeric     13       73       40   0.101        0        0
## 8       origin numeric      3        1      249   0.626        0        0
\end{verbatim}

Warning -\textgreater{} becasue some of the character strings aren't
properly structures integrers and so can't be translated to numeric

\hypertarget{split-data}{%
\section{5. Split data}\label{split-data}}

\begin{Shaded}
\begin{Highlighting}[]
\FunctionTok{set.seed}\NormalTok{(}\DecValTok{123}\NormalTok{)}
\NormalTok{n }\OtherTok{\textless{}{-}} \FunctionTok{nrow}\NormalTok{(autompg)}
\NormalTok{train }\OtherTok{\textless{}{-}} \FunctionTok{sample}\NormalTok{(n, }\FunctionTok{round}\NormalTok{(}\FloatTok{0.7} \SpecialCharTok{*}\NormalTok{ n))}
\NormalTok{autoTrain }\OtherTok{\textless{}{-}}\NormalTok{ autompg[train,]}
\NormalTok{autoTest }\OtherTok{\textless{}{-}}\NormalTok{ autompg[}\SpecialCharTok{{-}}\NormalTok{train,]}
\end{Highlighting}
\end{Shaded}

\hypertarget{fit-the-model}{%
\section{6. Fit the Model}\label{fit-the-model}}

\hypertarget{a.-linear-regression-models}{%
\subsection{6A. Linear Regression
Models}\label{a.-linear-regression-models}}

\begin{Shaded}
\begin{Highlighting}[]
\NormalTok{LRModel1 }\OtherTok{\textless{}{-}} \FunctionTok{lm}\NormalTok{(mpg }\SpecialCharTok{\textasciitilde{}}\NormalTok{., autoTrain) }\CommentTok{\# Model 1}
\CommentTok{\#summary(LRModel1)}

\NormalTok{LRModel2 }\OtherTok{\textless{}{-}} \FunctionTok{lm}\NormalTok{(mpg }\SpecialCharTok{\textasciitilde{}}\NormalTok{ weight }\SpecialCharTok{+}\NormalTok{ model\_year }\SpecialCharTok{+}\NormalTok{ origin, autoTrain) }\CommentTok{\# Model2}
\end{Highlighting}
\end{Shaded}

\hypertarget{b.-decission-tree-models}{%
\subsection{6B. Decission Tree Models}\label{b.-decission-tree-models}}

\begin{Shaded}
\begin{Highlighting}[]
\NormalTok{DTModel1 }\OtherTok{\textless{}{-}} \FunctionTok{rpart}\NormalTok{(mpg }\SpecialCharTok{\textasciitilde{}}\NormalTok{ ., autoTrain) }\CommentTok{\# Model 1}
\CommentTok{\#DTModel1$variable.importance}
\NormalTok{DTModel2 }\OtherTok{\textless{}{-}} \FunctionTok{rpart}\NormalTok{(mpg }\SpecialCharTok{\textasciitilde{}}\NormalTok{ displacement }\SpecialCharTok{+}\NormalTok{ weight }\SpecialCharTok{+}\NormalTok{ cylinders }\SpecialCharTok{+}\NormalTok{ horsepower, autoTrain) }\CommentTok{\# Model 2}
\end{Highlighting}
\end{Shaded}

Getting top 4 for the 2nd Model

\begin{Shaded}
\begin{Highlighting}[]
\FunctionTok{rpart.plot}\NormalTok{(DTModel1)}
\end{Highlighting}
\end{Shaded}

\includegraphics{Regression_files/figure-latex/unnamed-chunk-8-1.pdf}

\begin{Shaded}
\begin{Highlighting}[]
\FunctionTok{rpart.plot}\NormalTok{(DTModel2)}
\end{Highlighting}
\end{Shaded}

\includegraphics{Regression_files/figure-latex/unnamed-chunk-8-2.pdf}

\hypertarget{model-validation}{%
\section{7. Model Validation}\label{model-validation}}

\begin{Shaded}
\begin{Highlighting}[]
\NormalTok{predLR1 }\OtherTok{\textless{}{-}} \FunctionTok{predict}\NormalTok{(LRModel1, autoTest)}
\NormalTok{predLR2 }\OtherTok{\textless{}{-}} \FunctionTok{predict}\NormalTok{(LRModel2, autoTest)}
\NormalTok{predDT1 }\OtherTok{\textless{}{-}} \FunctionTok{predict}\NormalTok{(DTModel1, autoTest)}
\NormalTok{predDT2 }\OtherTok{\textless{}{-}} \FunctionTok{predict}\NormalTok{(DTModel2, autoTest)}
\end{Highlighting}
\end{Shaded}

\hypertarget{model-performance}{%
\section{8. Model Performance}\label{model-performance}}

\begin{Shaded}
\begin{Highlighting}[]
\CommentTok{\# Calculating RMSE}
\NormalTok{valLR1 }\OtherTok{\textless{}{-}} \FunctionTok{RMSE}\NormalTok{(predLR1, autoTest}\SpecialCharTok{$}\NormalTok{mpg)}
\NormalTok{valLR2 }\OtherTok{\textless{}{-}} \FunctionTok{RMSE}\NormalTok{(predLR2, autoTest}\SpecialCharTok{$}\NormalTok{mpg)}
\NormalTok{valDT1 }\OtherTok{\textless{}{-}} \FunctionTok{RMSE}\NormalTok{(predDT1, autoTest}\SpecialCharTok{$}\NormalTok{mpg)}
\NormalTok{valDT2 }\OtherTok{\textless{}{-}} \FunctionTok{RMSE}\NormalTok{(predDT2, autoTest}\SpecialCharTok{$}\NormalTok{mpg)}

\FunctionTok{sprintf}\NormalTok{(}\StringTok{"RMSE of Linear Model 1 is \%s"}\NormalTok{, valLR1)}
\end{Highlighting}
\end{Shaded}

\begin{verbatim}
## [1] "RMSE of Linear Model 1 is 3.2879450560213"
\end{verbatim}

\begin{Shaded}
\begin{Highlighting}[]
\FunctionTok{sprintf}\NormalTok{(}\StringTok{"RMSE of Linear Model 2 is \%s"}\NormalTok{, valLR2)}
\end{Highlighting}
\end{Shaded}

\begin{verbatim}
## [1] "RMSE of Linear Model 2 is 3.20023167178661"
\end{verbatim}

\begin{Shaded}
\begin{Highlighting}[]
\FunctionTok{sprintf}\NormalTok{(}\StringTok{"RMSE of Decission Tree Model 1 is \%s"}\NormalTok{, valDT1)}
\end{Highlighting}
\end{Shaded}

\begin{verbatim}
## [1] "RMSE of Decission Tree Model 1 is 3.48706632413932"
\end{verbatim}

\begin{Shaded}
\begin{Highlighting}[]
\FunctionTok{sprintf}\NormalTok{(}\StringTok{"RMSE of Decission Tree Model 2 is \%s"}\NormalTok{, valDT2)}
\end{Highlighting}
\end{Shaded}

\begin{verbatim}
## [1] "RMSE of Decission Tree Model 2 is 4.26075125701374"
\end{verbatim}

\begin{Shaded}
\begin{Highlighting}[]
\CommentTok{\# Calculationg R2\_Score}
\NormalTok{valLR1 }\OtherTok{\textless{}{-}} \FunctionTok{R2\_Score}\NormalTok{(predLR1, autoTest}\SpecialCharTok{$}\NormalTok{mpg)}
\NormalTok{valLR2 }\OtherTok{\textless{}{-}} \FunctionTok{R2\_Score}\NormalTok{(predLR2, autoTest}\SpecialCharTok{$}\NormalTok{mpg)}
\NormalTok{valDT1 }\OtherTok{\textless{}{-}} \FunctionTok{R2\_Score}\NormalTok{(predDT1, autoTest}\SpecialCharTok{$}\NormalTok{mpg)}
\NormalTok{valDT2 }\OtherTok{\textless{}{-}} \FunctionTok{R2\_Score}\NormalTok{(predDT2, autoTest}\SpecialCharTok{$}\NormalTok{mpg)}

\FunctionTok{sprintf}\NormalTok{(}\StringTok{"R2\_Score of Linear Model 1 is \%s"}\NormalTok{, valLR1)}
\end{Highlighting}
\end{Shaded}

\begin{verbatim}
## [1] "R2_Score of Linear Model 1 is 0.808512245475505"
\end{verbatim}

\begin{Shaded}
\begin{Highlighting}[]
\FunctionTok{sprintf}\NormalTok{(}\StringTok{"R2\_Score of Linear Model 2 is \%s"}\NormalTok{, valLR2)}
\end{Highlighting}
\end{Shaded}

\begin{verbatim}
## [1] "R2_Score of Linear Model 2 is 0.818592707592808"
\end{verbatim}

\begin{Shaded}
\begin{Highlighting}[]
\FunctionTok{sprintf}\NormalTok{(}\StringTok{"R2\_Score of Decission Tree Model 1 is \%s"}\NormalTok{, valDT1)}
\end{Highlighting}
\end{Shaded}

\begin{verbatim}
## [1] "R2_Score of Decission Tree Model 1 is 0.784616555127777"
\end{verbatim}

\begin{Shaded}
\begin{Highlighting}[]
\FunctionTok{sprintf}\NormalTok{(}\StringTok{"R2\_Score of Decission Tree Model 2 is \%s"}\NormalTok{, valDT2)}
\end{Highlighting}
\end{Shaded}

\begin{verbatim}
## [1] "R2_Score of Decission Tree Model 2 is 0.678438337798266"
\end{verbatim}

\hypertarget{conclusion}{%
\section{9. Conclusion}\label{conclusion}}

We can conclude that the best model is the Linear Regression Model 2
because it has the least RMSE and the highest R2\_Score

\end{document}
